\chapter{Talking Back: Input and Variables in Python}

\section{Introduction: Real-Life Hook}

In real life, communication is two-way:

\begin{itemize}
    \item Chat apps: you send a message and your friend replies  
    \item Voice assistants: you ask a question, it answers  
    \item ATM: you enter PIN, machine responds  
\end{itemize}

Similarly, Python can \textbf{take input from the user} and store it in variables.

\section{Concept: Getting Input from Users}

\begin{definition}{Input Function}
\texttt{input()} is a Python function that allows the program to take information from the user.
\end{definition}

\subsection{Example: Asking for Your Name}

\begin{verbatim}
name = input("What is your name? ")
print("Hello,", name)
\end{verbatim}

If the user types \texttt{Rafi}, the output will be:

\begin{verbatim}
What is your name? Rafi
Hello, Rafi
\end{verbatim}

\section{Input and Variables Together}

Variables can store the user’s input for later use in the program:

\begin{verbatim}
age = input("How old are you? ")
print("You are", age, "years old")
\end{verbatim}

\subsection{Observation}

- Whatever the user types is stored as a \textbf{string} by default.  
- You can convert it to numbers using \texttt{int()} or \texttt{float()}.

\begin{verbatim}
age = int(input("How old are you? "))
print("Next year you will be", age + 1)
\end{verbatim}

\section{Mini-Project: Personal Quiz Program}

Create a small Python program that asks the user about their favorite color, food, and hobby, then prints a summary:

\begin{verbatim}
# Personal Quiz
color = input("What is your favorite color? ")
food = input("What is your favorite food? ")
hobby = input("What is your favorite hobby? ")

print("Your favorite color is", color)
print("Your favorite food is", food)
print("Your hobby is", hobby)
\end{verbatim}

\section{Step-by-Step Diagram: Input and Output}

\begin{center}
\begin{tikzpicture}[node distance=1.5cm]
\node (start) [draw, ellipse, fill=green!20] {Start Program};
\node (input) [draw, rectangle, fill=yellow!20, below of=start] {Get Input from User};
\node (store) [draw, rectangle, fill=orange!20, below of=input] {Store Input in Variables};
\node (process) [draw, rectangle, fill=orange!20, below of=store] {Optional Processing};
\node (output) [draw, rectangle, fill=green!20, below of=process] {Display Result};
\draw[->] (start) -- (input);
\draw[->] (input) -- (store);
\draw[->] (store) -- (process);
\draw[->] (process) -- (output);
\end{tikzpicture}
\end{center}

\section{Key Concepts}

\begin{itemize}
    \item \texttt{input()} allows the program to communicate with the user.
    \item User input is stored in variables for later use.
    \item Input from the user is treated as a string unless converted to a number.
    \item Combining input with variables makes programs interactive.
\end{itemize}

\section{End-of-Chapter Exercises}

\begin{exercise}{Very Short Questions}
\begin{enumerate}
    \item What does the \texttt{input()} function do?
    \item How do you store user input in a variable?
    \item Why is user input treated as a string by default?
\end{enumerate}
\end{exercise}

\begin{exercise}{Short Answer Questions}
\begin{enumerate}
    \item Write a Python statement to ask the user for their favorite book and store it in a variable.
    \item How would you convert input to a number?
    \item Explain why combining input with variables is important for interactive programs.
\end{enumerate}
\end{exercise}

\begin{exercise}{Activity-Based Questions}
\begin{enumerate}
    \item Create a program that asks the user for three personal details and prints a summary.
    \item Modify the personal quiz to calculate the user’s age next year.
    \item Ask a classmate to run your program and observe different outputs.
\end{enumerate}
\end{exercise}

\section{Teacher Notes}

\marginnote{
\textbf{Learning Objectives:} 
\begin{itemize}
    \item Students learn to take input from users using Python.
    \item Students store input in variables.
    \item Students understand string vs numeric input.
\end{itemize}

\textbf{Teaching Tips:} 
\begin{itemize}
    \item Demonstrate input in Python shell or IDE live.
    \item Encourage students to try different inputs.
\end{itemize}

\textbf{Assessment:} 
\begin{itemize}
    \item Check if programs correctly store and display user input.
    \item Observe student understanding of string and number input.
\end{itemize}
}